\section{Related work and its limitations}\label{sec:related}

The literature this work sits in divides into three strands, and the gap it addresses lies
between them. Section~\ref{sec:rel-class} covers classifiers built on cough and
respiratory audio, Section~\ref{sec:rel-eval} covers what has since been established
about how their performance was measured, and Section~\ref{sec:rel-denoise} covers
denoising as a preprocessing step and states what the present work supplies that none of
these provide.

\subsection{Classifiers on cough and respiratory audio}\label{sec:rel-class}

Interest in classifying respiratory audio grew sharply during the COVID-19 pandemic, and
the resulting body of work established the pipeline this paper studies. An early mobile
system passed three second cough recordings to a cloud classifier and returned a
screening result, reporting that it was a screening aid rather than a clinical grade
test~\cite{imran2020}. A convolutional model trained on recordings collected through a web
form reported $98.5$ per cent sensitivity at $94.2$ per cent specificity against official
test results~\cite{laguarta2020}. Working from the Coswara corpus together with a smaller
clinical collection, a comparison of seven classifiers reported area under the curve of
$0.98$ separating positive from healthy coughs, and additionally observed that positive
coughs were some $15$ to $20$ per cent shorter~\cite{pahar2021}.

The supporting corpora were assembled over the same period, and their design decisions
shape what can be asked of them. Coswara records each participant performing a prescribed
sequence of breathing, cough and voice tasks and carries a participant
identifier~\cite{coswara2020}. COUGHVID gathered some $34{,}000$ recordings through a web
form and, because no human could annotate them all, shipped the output of a cough
detector as a per recording score~\cite{coughvid2021}. A larger crowdsourced collection
assembled over $53{,}000$ samples from more than $36{,}000$ participants and published
benchmarks for symptom and condition prediction~\cite{xia2021}. Outside the respiratory
setting, general purpose audio event corpora provide human annotations against a shared
ontology~\cite{gemmeke2017}, and one such corpus had its evaluation split exhaustively
verified by hand~\cite{fsd50k2022}, which is what makes it usable as an independent test
of a detector trained elsewhere.

Within this strand, the specific pipeline of interest is the one that converts a
recording into a spectrogram image and hands it to a convolutional network pretrained on
photographs. That approach reported $98.18$ per cent accuracy on a mel spectrogram
representation of cough audio~\cite{zhou2021}, and variants of it supply the architectures
used here~\cite{he2016,sandler2018,huang2017}.

\subsection{How the reported performance was measured}\label{sec:rel-eval}

The second strand is the subsequent examination of those results, and it is the reason
this paper is organised around measurement rather than around a new architecture. A
systematic review of machine learning applied to COVID-19 imaging screened $2{,}212$
studies, retained $62$ after quality assessment, and concluded that none of the models it
examined was of potential clinical use, attributing this to methodological flaws and
underlying biases rather than to insufficient model
capacity~\cite{roberts2021}. The same pattern was then demonstrated directly on audio. In
a study of $67{,}842$ participants, classifiers that reached $0.846$ area under the curve
unadjusted fell to $0.619$ once confounders including self-reported symptoms were matched,
and were outperformed by predictions from self-reported symptoms
alone~\cite{coppock2024}. A matched analysis on crowdsourced respiratory audio reported
$0.71$ area under the curve for a model evaluated on participant independent partitions,
against up to $0.90$ for the same material under a less careful evaluation
strategy~\cite{han2022}.

Three specific weaknesses recur, and each is addressed by a part of the present design.
The first is the split. Where a partition is made over recordings rather than
contributors, a contributor can appear on both sides of it. One study reports both
variants and so quantifies the cost: $98.18$ per cent under a random split against
$95.18$ per cent when test coughs came from recordings unseen in training, a gap it
presents without comment~\cite{zhou2021}. The second is what the learned representation
actually encodes. A recent cross-dataset evaluation of cough models found that
performance moderate within a corpus falls below $0.6$ area under the curve on
independent data, and that the representations cluster by recording device and corpus
rather than by the condition being detected~\cite{zhang2026}; two of its authors are also
authors of COUGHVID~\cite{coughvid2021}. The third is the absence of any estimate of run
to run variation, which leaves differences of well under one accuracy point reported as
findings with nothing to read them against.

A fourth weakness is specific to corpora whose labels come from a classifier. Where the
label is a thresholded detector score, a model is measured on its agreement with that
detector and not on whether a cough is present, and this distinction is frequently left
implicit. It cannot be resolved by better modelling, only by stating the regime and, where
possible, comparing across regimes.

\subsection{Denoising as a preprocessing step}\label{sec:rel-denoise}

The third strand is the preprocessing itself. Wavelet shrinkage has a well established
theory: the universal threshold with soft thresholding comes with an oracle
inequality~\cite{donoho1994}, a level adaptive threshold derived in a Bayesian framework
improves on it for image data~\cite{chang2000}, and an empirical Bayes rule with a mixture
prior estimates the sparsity of each scale from the data rather than assuming
it~\cite{johnstone2005}. All three rest on the same multiresolution
decomposition~\cite{mallat1989}, and the appeal for cough audio is that shrinkage
suppresses broadband noise while preserving transients, which is what a cough is.

What is missing is a comparison that isolates the effect. Where denoising is reported to
help, the denoised and undenoised conditions commonly differ in more than the denoising,
the shrinkage rule is often not named even though the three rules above behave quite
differently, and the comparison is typically between two single runs with no paired test
and no measure of seed variation. The work this paper extends illustrates all three: its
baseline and denoised conditions differ in architecture, in pretraining, in frequency axis
and in sample rate simultaneously, so no cause can be assigned to the difference between
them; its text describes soft thresholding while the routine it calls applies an empirical
Bayes posterior median by default; and its conclusion rests on one run per condition.

The present work therefore supplies what this configuration of strands leaves absent: a
denoising comparison under a named shrinkage rule, with the architecture, the frequency
axis, the partition and the training protocol held fixed, on a speaker disjoint partition
in which both classes come from the same participant and session, tested on paired per
recording decisions with a correction for multiple comparisons, repeated across seeds so
that the size of an effect can be read against the variation between runs, and carried
out across three label regimes rather than one. The contribution is the measurement, not
a new architecture.
